\section*{Political Violence, Networks, \& Civilian Populations: An Inclusive Approach for Conflict Evolution}

Our study of civil conflict takes, as its starting point, the belief that dynamics in intrastate conflict take place within a system in which each actor observes and reacts to the actions of others. Three primary features characterize complex systems of conflict: (1) dependencies between the actions of actors in the system; (2) the impact of armed groups on violence as they enter the conflict network; and (3) the ability of civilians to influence the strategic interactions of armed groups. While we view each factor as important and distinct, they are also deeply interrelated: the entrance of groups shapes dependencies in the conflict network, and these dependencies condition the ability of civilians to influence armed groups. We expand on each of these propositions below. 

\subsection*{Network Patterns in an Intrastate Conflict Context}

In the field of interstate conflict studies, network approaches are gaining more attention.\footnote{For example, see \citet{kinne:2012,metternich:etal:2015,minhas:etal:2016}.} These approaches have seen little usage in the study of civil wars even though the dynamics underlying many modern day intrastate conflicts involve interdependent interactions between multiple groups. The reason for studying dyadic interactions using a network-based approach is to acknowledge the possibility that the actions of actors in a system are contingent on one another. A significant amount of work has shown that by moving away from the assumption that interactions within a system are independent,  we can better account for the underlying process that generates relational data and make more accurate inferences.\footnote{For example, see \citet{snijders:1996,hoff:ward:2004,dorff:ward:2013,erikson:pinto:2014}.}  Three common forms of dependencies drive the creation of networked relationships: actor-level effects, second-order effects such as reciprocity, and third-order effects such as clustering among similar actors. We consider each of these below.

In social systems we often observe heterogeneity in actor centrality, revealing that some actors are more likely to initiate or receive ties (in this case, battles). Heterogeneity across actor centrality provides us with information on which set of events led to the creation of the observed network \citep{barabasi:reka:1999}.  Heterogeneity \textit{across} actors’ activity in a network produces homogeneous ties \textit{between} actors. For example, a particularly well-armed rebel group, $i$ might frequently target less resourced groups ($j$,$k$,$l$). Group resources, then, is a latent, unobserved attribute that accounts for why one rebel group is more likely to target a set of opposition groups. Thus, the relationships between the set of dyads $\{i \rightarrow j,\, i \rightarrow k,\, i \rightarrow l\}$ will be more similar to each other than they are to other interactions in the system, such as $m \rightarrow n$, because they all involve actor $i$ as the sender of the event \citep{kenny:lavoie:1984}.
\textcolor{red}{We can delete my added ex above if needed (since we have one below) but I really think this stuff was getting lost/confusing and it was useful to put it in more tangible terms at the onset of the explanation. I also understand if resources are too iffy.}

A straightforward example of why this type of dependence emerges in intrastate conflict networks involves government actors. Government actors likely have interests in asserting their control across the country. Specifically, in a scenario where rebel groups $\{j,k,l\}$ operate in different parts of a country, we would expect government actor $i$ to be more likely to initiate conflict with those rebel groups than an actor with non-competing interests. This logic naturally extends to non-government groups. For example, a rebel group may be acting upon ambitions of establishing control over broad swaths of the country, and, as a result, become a central initiator of conflicts in the network. For similar reasons we may observe heterogeneity in how likely actors are to receive conflict events, and, relatedly, find that actors initiating more battles are also likely to appear on the receiving end of violent attacks. Each of these processes creates dependence across observations that involve the same actor. 

An additional form of network dependence includes reciprocity, which is the notion that the interaction between $i \leftarrow j$ and $j \rightarrow i$ are dependent upon one another. The concept of reciprocity has deep roots in the study of relations between states \citep{richardson:1960,keohane:1989}, and it is no less relevant when considering how actors in an intrastate conflict context may act towards one another. \textcolor{red}{Needs more here, another walk through on rebels? Or perhaps just a bit more logic.}

Apart from these nodal and dyadic patterns that commonly arise in social systems, networks often exhibit patterns of dependence involving multiple actors. These types of patterns can broadly be summarized in terms of the concepts of stochastic equivalence and homophily \citep{wasserman:faust:1994}. Stochastic equivalence may manifest in a network in which actors cluster into groups, and the group determines an actor's patterns of relations with other groups. More concretely, a pair of actors $ij$ are stochastically equivalent if the probability of $i$ relating to, and being related to, by every other actor is the same as the probability for $j$ \citep{anderson:etal:1992}. Consider a case in which actors nest in groups based on, for example, a shared but unobserved ideology: $\mathbf{a}=\{i,j\}$ and $\mathbf{b}=\{k,l\}$. Here we may likely find that the relations of $i$ towards $k$ and $l$ is similar to how $j$ behaves towards those actors, and this induces a dependence in the relations of actors within the system. Specifically, the behavior of actors within the conflict system may be a function of a latent community to which they belong. Homophily is an additional type of dependence pattern involving multiple actors, and is typically used to explain the emergence of triads, cases in which actors $\{i,j,k\}$ are each connected and form a triangle, within social systems \citep{shalizi:thomas:2011}. In the intrastate setting, conflict triads may emerge when multiple actors are competing for control over the same territory or resource.

The principal idea underlying each of the dependence patterns discussed above is that the actions, or inactions, of any actor, are taking place within an interdependent system. Given this, the entrance or even exit of an actor may have a notable impact on the type of conflict patterns we observe across the network. We will demonstrate the importance of considering these effects in the context of the Nigeria conflict network. 

\subsection*{Entrance of Key Players}

An often undiscussed concept in networks and conflicts relates to the ripple effects that the entrance of a key player can have on a system. The entry of a particularly active or violent group could lead to an increase in conflict through a number of mechanisms. As \citet[p. 168]{kathman:wood:2015} have argued ``conflict systems are fluid, and competition varies in response to the arrival or exit of violent combat groups...", thus the entrance of certain groups increases the level of conflict. One possible mechanism through which violence might increase is through a process of rebel group outbidding, where in order for different groups to get support and resources from the population, they need to prove their conviction and capabilities through increasingly violent acts \citep{crenshaw:1981, bloom:2004, nemeth:2014}. The entrance of additional groups might also make ending conflict harder, as \citet{cunningham:etal:2009} notes, because the entrance of an additional (dominant) group makes de-escalating conflict more difficult as it increases the number of veto players in peace negotiations.

We also argue that the entrance of key players affects violence because of their strategic relationship with the government. The arrival of a new challenger group that successfully challenges the government may shift perceptions of the government as weak or vulnerable, leading other groups to more willingly challenge the government, as shown in \citet{walter:2006}. This particular dynamic also gives the government incentives to react with violence, rather than negotiation, to deter the development of future challengers. \citet[p. 613]{fjelde:nilsson:2012} find ``states with weak coercive power create opportunities for nonstate actors to engage in armed struggle against each other.'' 

Boko Haram is a perfect example of this effect on the Nigerian conflict network. Boko Haram gained such success that in mid-2014 they were able to effectively control swathes of territory in and around their home state of Borneo. The rise of Boko Haram and the challenges it posed to Nigeria's security apparatus provides incentives for other groups to increase their level of violence, even in dyads that do not involve Boko Haram.

\subsection*{Civilian Victimization in Armed Conflict}

%More often than not, analyses of armed conflict, civil war, and widespread criminal violence shine a spotlight on those actors that seek to wield power through collective violence against a given population or government. However, this narrative has begun to change. With uprisings across the last five years--the Tunisian revolution, the Egyptian revolution, the Syrian Civil War, and the 2014 protests in Venezuela--attention has shifted to recognize the potential power of the population to influence environments of extreme violence and repression. 

Compared to the number of studies focusing on economic \citep{collier:hoeffler:2004}, political \citep{humphreys:weinstein:2008}, and identity-based \citep{cederman:etal:2010a} drivers of conflict, existing research has largely neglected the role of civilians in influencing the behavior of armed actors. In particular we believe that civilians can influence combat dynamics in two major ways: through their responses to civilian victimization, and through their mobilization against violent actors. While the relationship between these factors and violence is important in their own right, we believe that by accounting for civilians, we can also create a model that better predicts the occurrence and evolution of conflict. However, the effect of civilians is not distinct from the effect of networks and key players, civilians' ability to influence dynamics of conflict are conditioned on the network structures of conflict.

While there is a robust debate over what causes civilian victimization,\footnote{For example, see \citet{valentino:etal:2004, downes:2006, kalyvas:2006, salehyan:etal:2014, prorok:appel:2014, humphreys:weinstein:2006}.} discussion of the consequences of civilian victimization, particularly as they relate to the potential for future conflict has been more limited \citep{hultman:2007, raleigh:2012}. %\citet{hultman:2007} finds that rebel group violence against civilians is a military strategy aimed at pressuring governments to concede to rebel challengers. 
\citet{berman:matanock:2015} argues that when rebels groups [governments] kill civilians, other civilians are more [less] likely to share information with the government: information allows the government to carry out attacks against insurgents, and the lack of information makes insurgent attacks less likely. \citet{condra:shapiro:2012} find support for this relationship in Iraq. Areas where US-led coalition forces kill civilians are associated with a higher likelihood of future insurgent attacks, and areas where insurgent forces kill civilians are associated with a lower risk of future insurgent attacks (against the US-led coalition).\footnote{\citet{lyall:2009}, however,  a contrary effect and demonstrates that indiscriminate violence against civilians in Chechnya is associated with a lower risk of rebel attacks, possibly because these attacks deprive insurgents of resources. }

These accounts of civil conflict primarily focus on a rebel-government dyad.\footnote{An exception is a 2015 study by \citeauthor{kathman:wood:2015}. They find that due to competition over resources the entry of additional rebel groups leads to an increase in civilian victimization, and that direct conflict between rebel groups increases civilian victimization.} In this dyad, conflict is zero-sum: conditions that enable the government to attack rebels make the rebels less able to attack the government, and vice versa.  If we follow this logic, then the dynamics of civilian victimization are conditioned by the network structure \citep{raleigh:2012}.\footnote{For an analogous discussion of how moving from dyadic to three player interactions destabilizes interstate conflict, see \citet{gallop:2017}.} If a civilian responds to violence by sharing information, the implications will depend on the network structure -- if they share information with a group that is a rival of the perpetrator (for example due to territorial competition and homophily) then civilian victimization will have a markedly different effect than if civilians share information with a group that is unrelated to the perpetrator, or even one that has common foes (a group that is stochastically equivalent to the perpetrator of civilian violence). Thus, while we believe that accounting for civilian victimization can us to better predict battles, for it to do so, we need to look at it in a context that accounts for the network structure.  %According to previous research on dyadic conflicts, civilians know the actors to rely on to obtain protection and revenge--the group that did not victimize them. %Typically, scholars conceptualize this group as the government.  

%Yet, in a networked conflict, the choice to respond to victimization by informing on the perpetrators (or siding with their opponents) is less clear. In this setting, civilians observe that the very existence of multiple challenger groups weakens the government. Thus in a multiparty conflict, the government is less likely to be viewed as capable of effective civilian protection when fighting between rebel groups increase \citep{fjelde:nilsson:2012}. In addition, if a rebel group preys on civilians, civilians may not be able to safely identify or contact the victimizers' rival rebel group \textit{nor} trust a weakened or uninterested government for protection.

%Whereas in dyadic conflicts we might expect civilian victimization to mitigate a group's ability to commit violent actions, in a multi-actor setting we do not expect civilian victimization to be associated with an armed group carrying out fewer attacks. In fact we expect the opposite to be the case because the incentives for groups to target civilians increase. In a multi-actor environment the countervailing effects of targeting civilians are weakened as civilians become less likely to report predation.  Groups might target civilians in order to consolidate resources or signal military strength \citep{kathman:wood:2015,hultman:2007}. In either case targeting civilians helps a group carry out attacks in the future. At the same time, we also expect groups who target civilians will become the victims of future attacks. First, because of retaliation or due to direct responses by challenger groups to aggressive behavior. Second, although the networked nature of conflict reduces civilians incentives to provide aid to rivals, it does not eliminate them.

\subsection*{Civilian Mobilization in Armed Conflict}

%Relatedly, operating in a conflict environment with many actors has implications for the effects of civilian mobilization. A large body of research already explores the origins of collective action and mobilization \citep{gurr:1970, opp:1988, tarrow:1994, tucker:2007}, and recent scholarship has assessed the effectiveness of ``maximalist campaigns" to show that countries are more likely to be democratic following nonviolent campaigns \citep{chenoweth:stephan:2011}.  Importantly, \citeauthor{celestino:gleditsch:2013} find that nonviolent campaigns destabilize regimes, but the trajectory of peace and democracy following nonviolent campaigns is conditional on the precise actions employed by campaign organizers. 

Building on the research agenda motivated by macro-level studies \citep{gurr:1970, opp:1988, tarrow:1994, tucker:2007,chenoweth:stephan:2011}, we investigate the link between civilian mobilization and violence between armed actors at the local level.  However, as articulated by \citet{celestino:gleditsch:2013}, even the macro-level relationship between nonviolent campaigns and state-level outcomes such as democracy or regime transitions remains unclear. Several possible logics link civilian mobilization and violence between armed groups.  As with victimization, we believe the network structure will condition which of these logics should apply. 

First, participants might achieve their stated aims and decrease violent conflict in their region. If protests successfully demand that the government invest greater resources in the prevention of violence, then we would expect fewer incidents of violence following protests. If protests also demand that non-state actors stop warfare with the state, this manifests as a public disapproval of attempts by the non-state group to coerce or control the civilian population. With enough discontent, protests disrupt the ability of non-state actors to operate in their region. This might be especially likely to occur in cases where conflict is driven by homophily -- if groups have similar population and resource bases, they should be especially sensitive to civilian mobilization because of the risk that it could drive the civilian populace to their rivals.

Alternatively, protests could lead to an increase in violence. As violence between armed actors (state and non-state) increases, non-state actors often victimize civilians to highlight the government's inability to protect the population. Even if protestors demand a non-violent solution to warfare, a plausible logic connecting protests and violence is that protests raise the cost of government inaction so that armed government actors are more likely to fight hard in the next round of contestation. This logic would suggest an antagonistic government actor, one who is politically motivated to demonstrate competency through violent policy. Networks might also play a role here because protests against the government will have a different tenor depending on the shape of the conflict network. In cases where multiple non-state actors are fighting amongst each other (potentially due to homophily) protests might encourage the government to try to suppress violence, but in cases where there is a lot of stochastic equivalence, because many groups are fighting with the government, or cases where much of the governments violence is reciprocal, we might expect to see the government respond to protests by increasing their levels of violence.

%Finally, there might be no relationship between civilian protests and violence. Protests might have other benefits and thus persist for reasons outside of the stated goals driving protest organizers. For example, non-violent activism via protests might stimulate community between survivors of violence and foster a sense of purpose and belonging in the midst of crisis. Additionally, protests against insecurity might target at both types of armed actors (i.e., the state and the non-state challenger), these actors could respond to protests through other means other than violence. For instance, non-state armed groups could make public appeals to civilians, saying that they are there for civilians' protection. Similarly, governments might respond with media campaigns, speeches, or through a general effort to try to divert political attention from violent events. A second mechanism driving a null relationship between protests and violent events is participant risk. Following a protest, participants could be targeted because of their activism, possibly resulting in fewer protests over time but with no clear consequences for violence between armed groups.

In this paper, we test these competing logics using the case of Nigeria. In particular, we are interested in how conflict evolves following certain strategic actions: civil mobilization against military actors and civilian victimization by those actors. We investigate whether riots and protests targeting a group make them more (or less) likely to both attack other groups and be attacked by groups in the future. Similarly, we look at whether groups which target civilians are involved in more attacks, and are attacked more in the future.

The contribution of this study is two-fold. We explain and demonstrate the importance of considering conflict processes through a network framework. In doing so, we suggest that networks of conflict constitute a meaningful outcome of interest and depart from the vast majority of the literature which focuses on dyadic outcomes to measure conflict intensity or duration. Second, we contribute to a growing scholarship concerned with the influence of civilian victimization and mobilization on trajectories of violence and stability. In particular, we argue that combining these focuses are necessary because many of the effects of the civilian population are tied to the structure of the conflict network and the dependencies therein. 
